\begin{helper}
Fix \(I\subseteq[n]\), the integers \(k,\ell_R,\ell_B\), and the parameters
\(m,p,\varepsilon,\varepsilon_1,\varepsilon_2\), with \(0<p<1\).  Let
\[
S=\{\omega:
\noderef{TwoBiteProjectionSizeEvent}(I,k,\ell_R,\ell_B,\omega)\}.
\]
There is a finite partition \(S=\dot\bigcup_{h\in H}C_h\) into exposure
histories with the following properties.  For each history \(h\), there are a
recorded coordinate set \(R_h\), a terminal candidate coordinate set \(U_h\),
a terminal order \(L_h\) listing \(U_h\), and a representative
\(\omega_h\in C_h\), such that \(C_h\) is exactly the cylinder fixing the
injection of \(\omega_h\) and the truth table on \(R_h\).  For every
\(\omega\in C_h\), the recorded set is exactly the pre-terminal recorded
set:
\[
e\in R_h
\quad\Longleftrightarrow\quad
e\in\noderef{TwoBitePreTerminalRecordedPairs}(\omega,\varepsilon,I).
\tag{0}
\]
Every coordinate in \(R_h\) is represented in canonical increasing endpoint
order.

The terminal set is the full unrecorded terminal coordinate universe:
\[
e\in U_h
\quad\Longleftrightarrow\quad
e\in\noderef{TwoBiteTerminalCoordinateUniverse}(m)
\hbox{ and }e\notin R_h.
\tag{1}
\]
Thus every coordinate in \(U_h\) is outside \(R_h\), is a distinct non-loop
red or blue base-edge coordinate, and is represented with increasing endpoint
order.  Since \(R_h\) has the same canonical orientation, no element of
\(U_h\) can be indirectly fixed by a reversed coordinate in \(R_h\).  The set
\(U_h\) may be larger than the final staged-open set.
For every \(\omega\in C_h\),
\[
e\in\noderef{TwoBiteStagedOpenPairs}(\omega,\varepsilon,I)
\quad\Longrightarrow\quad
e\in U_h.
\tag{2}
\]
The set \(\noderef{TwoBiteStagedOpenPairs}(\omega,\varepsilon,I)\) is not
claimed to be fixed on \(C_h\); it is an adaptive subset of \(U_h\).

Conditional on \(C_h\), the coordinates in \(U_h\) have product
Bernoulli-\(p\) law.  Equivalently, for every subset \(A\subseteq U_h\),
\[
\Pr(\text{exactly the coordinates in }A\text{ are present on }U_h\mid C_h)
=p^{|A|}(1-p)^{|U_h|-|A|}.
\tag{3}
\]
The terminal order is prefix safe: when a coordinate is reached, whether it is
pre-terminal recorded, staged-open, touched, or \(C^+(I)\) same-color closed
is determined by the history and the earlier terminal-coordinate values, not
by its own Bernoulli value or by later coordinates.

Moreover, for every \(B\ge0\), any upper bound \(B\) for
\[
\Pr(I\text{ is independent in }\noderef{TwoBiteFinalGraph}
       \text{ and }
       \noderef{TwoBiteRegularityEvent}
       (k,\varepsilon,\varepsilon_1,\varepsilon_2)
       \mid C_h)
\]
which is uniform over all histories \(h\) also holds after conditioning only
on \(S\).
\end{helper}
\begin{proof}
Apply \(\noderef{FixedSetExposureHistoryCylinder}\) to the fixed data
\(I,k,\ell_R,\ell_B,m,p,\varepsilon\).  It gives a finite type \(H\), a
history predicate \(C_h\), recorded sets \(R_h\), terminal sets \(U_h\),
orders \(L_h\), and representatives \(\omega_h\).  In the Lean statement of
the present node we take
\[
\mathrm{hist}(h,\omega)\quad\text{to mean}\quad \omega\in C_h,
\]
and we use the same \(R_h,U_h,L_h,\omega_h\) as the witnesses for the
history \(h\).

The history-cylinder child supplies the exact partition data needed for the
first two conjuncts:
\[
\omega\in S \Longleftrightarrow \exists h\in H,\ \omega\in C_h,
\tag{4}
\]
and
\[
h\ne h'\Longrightarrow C_h\cap C_{h'}=\varnothing. \tag{5}
\]
For each \(h\), it also supplies
\[
\omega_h\in C_h, \tag{6}
\]
the exact cylinder identity
\[
\omega\in C_h
\Longleftrightarrow
\omega\hbox{ has the same injection as }\omega_h
\hbox{ and agrees with }\omega_h\hbox{ on }R_h, \tag{7}
\]
the exact recorded-set identity
\[
\omega\in C_h
\Longrightarrow
\bigl(e\in R_h\Longleftrightarrow
e\in\noderef{TwoBitePreTerminalRecordedPairs}(\omega,\varepsilon,I)\bigr),
\tag{7a}
\]
the recorded-coordinate hygiene
\[
e\in R_h\Longrightarrow e\hbox{ is written with increasing endpoints},
\tag{7'}
\]
the order facts
\[
L_h.\mathrm{Nodup}
\quad\text{and}\quad
L_h.\mathrm{toFinset}=U_h, \tag{8}
\]
the terminal-universe identity
\[
e\in U_h
\Longleftrightarrow
e\in\noderef{TwoBiteTerminalCoordinateUniverse}(m)
\hbox{ and }e\notin R_h, \tag{9}
\]
the disjointness of terminal and recorded coordinates
\[
e\in U_h\Longrightarrow e\notin R_h, \tag{10}
\]
and the increasing non-loop coordinate hygiene
\[
e\in U_h\Longrightarrow
\begin{cases}
a<b,& e=\operatorname{red}(a,b),\\
a<b,& e=\operatorname{blue}(a,b).
\end{cases}
\tag{11}
\]
It further gives the staged-open containment
\[
\omega\in C_h,\quad
e\in\noderef{TwoBiteStagedOpenPairs}(\omega,\varepsilon,I)
\Longrightarrow e\in U_h, \tag{12}
\]
and the prefix-safety clause: if \(L_h=P,e,Q\), then agreement on the
injection, on \(R_h\), and on \(P\) preserves, at the current coordinate
\(e\), the truth values of pre-terminal recordedness, staged-openness,
touchedness, and \(C^+(I)\) same-color closedness.  Thus all non-product
structural conjuncts of the present statement are copied directly from
\(\noderef{FixedSetExposureHistoryCylinder}\).

Now fix one history \(h\).  Apply
\(\noderef{FixedSetCylinderProductLaw}\) with the history event \(C_h\), the
recorded set \(R_h\), the product coordinate set \(U_h\), and the
representative \(\omega_h\).  In the notation of the Lean statement of that
child, this is the instantiation
\[
\mathrm{hist}:=C_h,\qquad
\mathrm{recorded}:=R_h,\qquad
\mathrm{product}:=U_h,\qquad
\omega_0:=\omega_h .
\]
Its hypotheses are exactly available.  The inequalities \(0<p\) and \(p<1\)
are the two standing Bernoulli-range assumptions of the present node.  The
representative condition (6) gives \(C_h(\omega_h)\).  The equivalence (7) is
the child theorem's required cylinder description: a sample is in \(C_h\)
precisely when it has the representative injection and agrees with
\(\omega_h\) on all recorded coordinates \(R_h\).  The terminal-recorded
disjointness (10) gives the hypothesis
\[
\forall e\in U_h,\ e\notin R_h,
\]
the recorded-coordinate hygiene (7') gives the product law's
canonical-recorded-coordinate hypothesis, and the increasing non-loop
coordinate hygiene (11) gives the final coordinate-shape hypothesis of
\(\noderef{FixedSetCylinderProductLaw}\).  Hence, for every
\(A\subseteq U_h\),
\[
\noderef{TwoBiteConditionalProbability}
\left(
\{\omega:\forall e\in U_h,\,
\noderef{TwoBiteEdgeCoordinateValue}(\omega,e)
\Longleftrightarrow e\in A\}
\mid C_h
\right)
=p^{|A|}(1-p)^{|U_h|-|A|}. \tag{13}
\]
This is exactly the product-law conjunct in the Lean statement.  The phrase
``\(e\) is present'' in (13) is interpreted through
\(\noderef{TwoBiteEdgeCoordinateValue}\): for a red tagged coordinate it is
membership of the corresponding edge in \(\noderef{TwoBiteRedGraph}\), and
for a blue tagged coordinate it is membership in
\(\noderef{TwoBiteBlueGraph}\).

The product law is deliberately a law on the fixed set \(U_h\), not a claim
that \(\noderef{TwoBiteStagedOpenPairs}(\omega,\varepsilon,I)\) is fixed on
the cell.  The only relation between the adaptive staged-open set and the
product set is the one-way containment (12), valid separately for each
\(\omega\in C_h\).  During the terminal reveal the coordinates of \(U_h\) are scanned
in the order \(L_h\).  If \(L_h=P,e,Q\), prefix safety says that after
conditioning on \(C_h\) and on the already exposed prefix values \(P\), the
pre-terminal recorded predicate and the three predicates deciding the paper's
terminal step for \(e\)--staged-open, touched, and same-color closed--are
already determined.  The Bernoulli value of \(e\) itself has not been used,
and no value from the suffix \(Q\) has been used.  Therefore the sequential
RISI argument may decide whether the current coordinate is eligible using only
the prefix information and then charge the fresh \(p\) or \(1-p\) factor for
the current coordinate from (13).  This is the needed adaptive replacement for the false fixed
staged-open-set assertion.

It remains to prove the finite averaging transfer.  Let
\[
\mathcal E(\omega)=
\bigl(I\text{ is independent in }\noderef{TwoBiteFinalGraph}(\omega)\bigr)
\land
\noderef{TwoBiteRegularityEvent}
      (k,\varepsilon,\varepsilon_1,\varepsilon_2,\omega).
\]
The regularity event is part of the numerator event \(\mathcal E\) only.  It
is not part of any conditioning event in this node: the final denominator is
the projection-size event \(S\), and the cell denominators are the history
events \(C_h\).

Assume \(B\ge0\) and that for every \(h\)
\[
\noderef{TwoBiteConditionalProbability}(n,m,p;\mathcal E\mid C_h)\le B.
\tag{14}
\]
If \(\Pr(S)=0\), then by the definition of
\(\noderef{TwoBiteConditionalProbability}\) the conditional probability
\(\Pr(\mathcal E\mid S)\) is \(0\), and the desired bound follows from
\(B\ge0\).  Hence assume \(\Pr(S)>0\).

The finite cover (4) and disjointness (5) imply that the events
\(\mathcal E\cap C_h\) are pairwise disjoint and have union
\(\mathcal E\cap S\): cover (4) places every sample of \(S\) in at least one
history cell, disjointness (5) makes that cell unique, and intersecting with
\(\mathcal E\) preserves both the union and the disjointness.  Since
\(\noderef{TwoBiteEventProbability}\) is a finite sum of nonnegative sample
weights,
\[
\Pr(\mathcal E\cap S)=\sum_{h\in H}\Pr(\mathcal E\cap C_h),
\qquad
\Pr(S)=\sum_{h\in H}\Pr(C_h). \tag{15}
\]
For an individual history \(h\), there are two cases.  If \(\Pr(C_h)=0\),
then \(\mathcal E\cap C_h\subseteq C_h\) and nonnegativity of the finite
sample-weight sum give
\[
\Pr(\mathcal E\cap C_h)=0\le B\Pr(C_h). \tag{16}
\]
If \(\Pr(C_h)>0\), then the nonzero branch of conditional probability and
(14) give
\[
\frac{\Pr(\mathcal E\cap C_h)}{\Pr(C_h)}
=\Pr(\mathcal E\mid C_h)\le B.
\]
Multiplying by the positive number \(\Pr(C_h)\) gives
\[
\Pr(\mathcal E\cap C_h)\le B\Pr(C_h). \tag{17}
\]
Thus (16) or (17) holds for every history \(h\).  Summing over the finite
history type, using \(\texttt{Finset.sum\_le\_sum}\) on the pointwise
inequalities, and then using (15),
\[
\Pr(\mathcal E\cap S)
\le
B\sum_{h\in H}\Pr(C_h)
=B\Pr(S).
\]
Dividing by the positive number \(\Pr(S)\) yields
\[
\Pr(\mathcal E\mid S)
=\frac{\Pr(\mathcal E\cap S)}{\Pr(S)}
\le B.
\]
This proves the RISI transfer from history-uniform bounds to conditioning
only on \(\noderef{TwoBiteProjectionSizeEvent}\), with regularity kept
solely in the numerator.
\end{proof}
